\documentclass[conference]{IEEEtran}
\IEEEoverridecommandlockouts
% The preceding line is only needed to identify funding in the first footnote. If that is unneeded, please comment it out.
\usepackage{cite}
\usepackage{amsmath,amssymb,amsfonts}
\usepackage{algorithmic}
\usepackage{graphicx}
\usepackage{multirow}
\usepackage{textcomp}
\usepackage{xcolor}
\usepackage{hyperref}
\usepackage{url}
\def\BibTeX{{\rm B\kern-.05em{\sc i\kern-.025em b}\kern-.08em
    T\kern-.1667em\lower.7ex\hbox{E}\kern-.125emX}}
\begin{document}

\title{OpenT2D Explorer: An Open-Source Clinical Analytics Dashboard for the AI-READI Dataset}

\author{\IEEEauthorblockN{Md Basit Azam$^{*}$,~\IEEEmembership{Graduate Student Member,~IEEE} and  Sarangthem Ibotombi Singh,~\IEEEmembership{Member,~IEEE} }
\IEEEauthorblockA{\\ Department of Computer Science \& Engineering, Tezpur University, Napaam - 784 028, Tezpur, Assam, INDIA.\\  \texttt{\{mdbasit, sis\}@tezu.ernet.in} \\
}
\thanks{$^*$ Corresponding Author: Azam, M.B  }
}

\maketitle

\begin{abstract}
The AI-READI Consortium Dataset v3.0.0 comprises 2,280 subjects with nine co-registered data modalities for type~2 diabetes (T2D) research, totalling 3.82~TB of heterogeneous data in formats such as waveform-database (WFDB) electrocardiogram (ECG) signals, Open~mHealth JSON continuous glucose monitoring (CGM) streams, and Observational Medical Outcomes Partnership (OMOP) common data model (CDM) clinical tables. Assessing whether such a dataset is \emph{ready} for predictive modelling is difficult precisely because the modalities are stored in incompatible formats with non-standard metadata. We present \emph{OpenT2D Explorer}, an open-source interactive dashboard (FastAPI backend, React frontend) that operationalises multimodal data-readiness assessment along three reusable axes through three purpose-built Exploratory Data Analysis (EDA) tools: (1)~a Temporal Overlap Visualizer (TOV) that anchors CGM and ECG windows to the clinical-visit date; (2)~a Multimodal Co-missingness Matrix quantifying pairwise modality file co-presence per study group; and (3)~a Signal Quality module reporting CGM dropout and Philips automatic ECG interpretation verdicts with cohort-filtering guidance. A distinguishing methodological contribution is a WFDB timestamp-recovery technique that extracts recording dates and interval measurements from non-standard free-text comment fields (100.0\% extraction success across N=2,251 records). We evaluate the system with a factual capability comparison against existing tools (consortium Jupyter notebooks, OHDSI~ATLAS, generic profilers), an author-estimated task-time analysis, and a four-part data-validity assessment. A representative five-step workflow shows that 2,220/2,280 participants (97.4\%) have concurrent ECG, CGM, and clinical files, and that HbA1c exhibits a monotonic clinical gradient (Healthy 5.5\%, Pre-DM 5.8\%, Oral~Med.\ 6.3\%, Insulin 7.0\%) matching the ADA treatment target.
\end{abstract}

\begin{IEEEkeywords}
Clinical dashboard, continuous glucose monitoring (CGM), electrocardiogram (ECG), multimodal diabetes data, data readiness, OMOP
\end{IEEEkeywords}

\section{Introduction}
The AI-READI project produced a large multimodal Type~2 Diabetes dataset and participant cohort for T2D research \cite{ai-readi_consortium_ai-readi_2024}. The v3.0.0 release \cite{ai-readi_consortium_flagship_2025} contains 2,280 participants collected at three clinical sites University of Washington (UW), University of California San Diego (UCSD), and University of Alabama at Birmingham (UAB) across nine co-registered modalities spanning the cardiac, retinal, wearable, and clinical categories.

Although the dataset is packaged following FAIR principles \cite{wilkinson_fair_2016}, using it for deep learning requires parallel parsing of WFDB \cite{moody_wfdb_nodate} binary ECG records, Open~mHealth JSON CGM files, and OMOP CDM tables \cite{hripcsak_george_observational_2015, noauthor_omop_nodate}. Before training a model, a researcher must answer three data-readiness questions: (1)~do the target modalities exist for enough participants? (2)~are the high-frequency streams captured in the same period relative to the clinic visit? and (3)~are the raw signals of sufficient quality to use without further curation? These three questions \emph{coverage}, \emph{temporal co-registration}, and \emph{signal integrity} are not specific to AI-READI; they recur in every effort to assemble a multimodal clinical training set, yet no single tool answers them for AI-READI today.

We present \emph{OpenT2D Explorer}, an open-source dashboard that answers 1--3 through three EDA tools. Rather than reporting new clinical findings, we contribute the \emph{methodology and infrastructure} that turn a raw, heterogeneously formatted release into a quality-verified, well-scoped cohort. Our contributions are:
\begin{itemize}
\item We formalise multimodal readiness as coverage $\times$ temporal co-registration $\times$ signal integrity, and instantiate each axis as an interactive tool; the framing is dataset-agnostic and transfers to other OMOP/WFDB/Open~mHealth cohorts.
\item AI-READI WFDB headers leave the standard \texttt{base\_date}/\texttt{base\_time} fields empty; we recover recording dates, firmware, filter settings, and interval measurements from non-standard free-text comment fields, enabling cross-modal temporal alignment that is otherwise impossible.
\item We report a capability comparison against existing tools (Table~\ref{tab:comparison}), an author-estimated task-time analysis (Table~\ref{tab:timesavings}), and a four-part validity assessment (Section~\ref{sec:validation}). The system is released under the MIT License.
\end{itemize}

\section{Related Work and Positioning}
\label{sec:related}
Three families of tools partially overlap with our problem. \emph{Consortium Jupyter notebooks} shipped with AI-READI demonstrate per-modality loading but leave cross-modal coverage, temporal alignment, and quality gating to be hand-coded by each researcher. \emph{OHDSI ATLAS} \cite{hripcsak_george_observational_2015} is a mature cohort-definition and characterisation tool for OMOP CDM data, but it is confined to the relational clinical tables and has no notion of WFDB waveforms or Open~mHealth CGM streams exactly the high-frequency modalities that dominate AI-READI. \emph{Generic data profilers} such as ydata-profiling \cite{ydata_profiling} and validation frameworks such as Great Expectations \cite{great_expectations} operate on flat tabular frames and are format-agnostic, so they cannot parse binary ECG headers or reconcile timestamps across modalities. External EHR resources such as MIMIC-IV \cite{johnson_mimic_2023} and eICU \cite{pollard_eicu_2018} motivate the OMOP-to-LOINC mapping we embed (Section~\ref{sec:engineering}) but are unrelated datasets.

Table~\ref{tab:comparison} contrasts these tools on the capabilities that matter for AI-READI readiness. The gap OpenT2D Explorer fills is the combination of (i)~multimodal \emph{file} co-presence, (ii)~binary-waveform and CGM parsing with quality metrics, and (iii)~cross-modal \emph{temporal} co-registration in one no-code, on-premises application.

\begin{table}[t]
    \centering
    \caption{Capability comparison with representative existing tools. \checkmark: supported; Part.: partial/manual; --: not supported.}
    \label{tab:comparison}
    \resizebox{\columnwidth}{!}{%
    \begin{tabular}{@{} l c c c c @{}}
        \hline
        \textbf{Capability} & \textbf{Jupyter} & \textbf{ATLAS} & \textbf{Profilers} & \textbf{Ours} \\
        \hline
        Multimodal file co-presence      & Part. & --    & --    & \checkmark \\
        WFDB ECG comment/interval parse  & Part. & --    & --    & \checkmark \\
        CGM JSON dropout quality         & Part. & --    & Part. & \checkmark \\
        OMOP clinical measurements       & --    & \checkmark & -- & \checkmark \\
        Cross-modal temporal alignment   & --    & --    & --    & \checkmark \\
        No-code interactive UI           & --    & \checkmark & \checkmark & \checkmark \\
        On-premises / DUA-compliant      & \checkmark & \checkmark & \checkmark & \checkmark \\
        AI-READI--specific               & Part. & --    & --    & \checkmark \\
        \hline
    \end{tabular}
    }
\end{table}

\section{System Design}

\subsection{Architecture}

\begin{figure*}
    \centering
    \includegraphics[width=\linewidth]{aireadi_architecture.pdf}
    \caption{System architecture: FastAPI backend with service modules and a REST API, React frontend with three EDA tool components. The dataset directory is mounted read-only.}
    \label{fig:aireadi_architecture}
\end{figure*}

OpenT2D Explorer follows a client--server architecture (Fig.~\ref{fig:aireadi_architecture}). A Python FastAPI \cite{noauthor_fastapi_nodate} backend serves a React~18 frontend via a REST API. All data access is read-only and fully on-premises; no participant data leaves the researcher's environment, satisfying the AI-READI data-use agreement (DUA). A Docker Compose deployment is provided (tested on a Google Cloud VM: e2-highmem-4, 4\,vCPU, 32\,GiB RAM). The backend exposes three router modules (\texttt{/api/cohort}, \texttt{/api/patients}, \texttt{/api/eda}) backed by seven format-specific service modules. A startup hook pre-loads the 2,280-row \texttt{participants.tsv} index into an LRU cache, keeping cohort-level responses under 50\,ms. The heaviest endpoint, \texttt{/api/eda/signal-quality} (scanning up to 200 CGM and 150 ECG files), completes in $\approx$18.2\,s cold and 9.0\,s warm, reflecting sequential filesystem reads across participant directories acceptable for an on-demand EDA endpoint. The three EDA tools (Section~\ref{sec:eda}) are served exclusively through \texttt{/api/eda}.

\subsection{Key Engineering Decisions}
\label{sec:engineering}
\textbf{ECG timestamp recovery.} AI-READI WFDB records do not populate the standard \texttt{base\_date}/\texttt{base\_time} header fields; recording dates are instead embedded as free-text comments (e.g., \texttt{validation\_date: 20241014}). Our \texttt{eda\_service.py} implements a structured comment parser extracting this timestamp alongside firmware version, filter settings (HP 0.15\,Hz, LP 100\,Hz, Notch 60\,Hz), HR/PR/QRS/QT/QTc intervals, and Philips machine-interpretation flags. This non-obvious step is what enables temporal alignment of ECG with CGM windows (contribution~C2). When \texttt{recording\_date} cannot be extracted it is set to \texttt{None} and the timeline displays ``---'' rather than raising an error, ensuring graceful degradation. The \texttt{has\_abnormal\_flag} field is set when an abnormal keyword appears in \emph{any} comment field, intentionally capturing subthreshold findings (e.g., minimal anterior ST elevation annotated alongside an ``OTHERWISE NORMAL ECG'' verdict) to give researchers the most sensitive flag for downstream exclusion.

\textbf{OMOP-to-LOINC mapping and unit validation.} The \texttt{measurement.csv} table \cite{noauthor_omop_nodate} contains 242,279 rows mapped to OMOP concept IDs. The dashboard embeds a static concept dictionary providing LOINC codes and canonical units for all 23 measured concepts, e.g., HbA1c (OMOP~3004410~$\to$~LOINC~4548-4, \%) and fasting glucose (OMOP~3004501~$\to$~LOINC~2345-7, mg/dL), enabling cross-reference with external EHR datasets such as MIMIC-IV \cite{johnson_mimic_2023} or eICU \cite{pollard_eicu_2018}. Unit homogeneity was verified via the \texttt{unit\_source\_value} column: all HbA1c rows use NGSP percent with no mmol/mol entries, and HbA1c, the lipid panel, creatinine, BUN, CRP, and MoCA were confirmed unit-homogeneous. A subset (anthropometrics, vitals, hematology) carries blank \texttt{unit\_source\_value}, and Insulin/C-peptide (\texttt{ng/L}) and Calcium (\texttt{mEq/L}) carry non-canonical units; these are handled by a physiologically plausible range filter (e.g., HbA1c 3--20\%, BMI 10--80\,kg/m$^2$) rather than unit conversion, since their values fall within expected ranges regardless of label. The service layer implements a conversion registry for future mixed-unit datasets and surfaces the excluded-row count in the API response for transparency ($<$0.2\% exclusion rate, consistent with isolated data-entry errors rather than systematic inconsistency).

\section{EDA Tools}
\label{sec:eda}

\subsection{Temporal Overlap Visualizer (TOV)}
\label{sec:temporal_overlap_visualizer}
Temporal alignment is the most critical preparation step for non-invasive glucose prediction from ECG. The TOV renders a per-participant Gantt-style timeline with three streams anchored to the clinical-visit date from \texttt{visit\_occurrence.csv}: (1)~the visit timestamp (point event); (2)~the CGM window as a continuous bar from first to last Open~mHealth timestamp; and (3)~the ECG date recovered from \texttt{.hea} comments. All comparisons operate at calendar-date granularity; sub-day alignment is not verifiable from the available metadata and is discussed in Section~\ref{sec:discussions}.

Figure~\ref{fig:temporal_overlap_visualizer_1023} shows Participant~1023. The visit occurred on 2023-08-30 at 08:55; CGM monitoring began the same calendar day at 17:53~UTC and ran for 9.9~days (2,816 readings, 1.4\% dropout). The recovered ECG \texttt{validation\_date} is 2024-10-14, a follow-up acquired $\approx$14~months after baseline, consistent with the AI-READI longitudinal protocol. This illustrates a key finding: while CGM co-registers with the baseline visit by design, ECG recordings are \emph{not} universally same-day and must be verified per participant using the TOV. The interface also surfaces automated ECG intervals (HR, QTc) and interpretation flags inline, letting a researcher exclude participants with clinically significant findings (e.g., QTc~$>$~500\,ms, arrhythmia) before building a training set.

\begin{figure*}[htbp]
    \centering
    \includegraphics[width=\linewidth]{Figure 2.png}
    \caption{TOV for Participant 1023. CGM monitoring (9.9 days, 2,816 readings) co-registers with the 2023-08-30 baseline visit. The ECG date (2024-10-14) reflects a longitudinal follow-up, illustrating the need for per-participant temporal verification.}
    \label{fig:temporal_overlap_visualizer_1023}
\end{figure*}

\subsection{Multimodal Co-missingness Matrix}
The Co-missingness Matrix reports how many participants have all required files for a given modality pair. A symmetric 9$\times$9 heatmap (Fig.~\ref{fig:co_missingness_matrix}) is constructed: cell~$(i,j)$ gives the number and percentage of participants with files from both modalities $i$ and $j$; the diagonal gives single-modality counts. The matrix captures file \emph{presence} (from the boolean modality columns in \texttt{participants.tsv}), not timestamp co-registration, which is confirmed individually with the TOV (Section~\ref{sec:temporal_overlap_visualizer}). This two-stage approach, file-level filtering followed by temporal-overlap inspection, is the intended usage. A study-group filter and a triple-overlap module reveal ECG, CGM, and clinical data jointly, which are vital for supervised glucose estimation. Table~\ref{tab:triple_overlap} reports results for all study groups.

\begin{figure}[htbp]
    \centering
    \includegraphics[width=\linewidth]{Figure 3.png}
    \caption{Multimodal Co-missingness Matrix for the full AI-READI cohort (N=2,280). Color intensity encodes co-presence percentage.}
    \label{fig:co_missingness_matrix}
\end{figure}

\begin{table}[htbp]
    \centering
    \caption{ECG + CGM + Clinical Data Triple Overlap by Study Group}
    \label{tab:triple_overlap}
    \begin{tabular}{@{} l c c c @{}}
        \hline
        \textbf{Study Group} & \textbf{N (Total)} & \textbf{N (Triple)} & \textbf{Coverage (\%)} \\
        \hline
        Healthy                 & 776   & 750   & 96.6 \\
        Pre-DM / Lifestyle      & 560   & 547   & 97.7 \\
        Oral / Non-insulin Med. & 686   & 675   & 98.4 \\
        Insulin Dependent       & 258   & 248   & 96.1 \\
        \hline
        \textbf{Total}          & \textbf{2,280} & \textbf{2,220} & \textbf{97.4} \\
        \hline
    \end{tabular}
\end{table}

\subsection{Signal Quality \& Sensor Provenance}
The Signal Quality module provides two complementary quality views for principled exclusion decisions. \textbf{CGM quality.} For a performance-optimised sample (the first 200 in \texttt{participants.tsv} ordering a deliberate trade-off for sub-second response), per-participant dropout is
\begin{equation}
    \frac{\text{expected} - \text{actual readings}}{\text{expected}} \times 100\%,
\end{equation}
with $\text{expected} = \text{days\_covered} \times 288$ (Dexcom~G6 5-minute interval \cite{dexcom_guides}). Dropout and monitoring-duration distributions are rendered as interactive histograms, so a researcher can read off what fraction of participants would survive, e.g., a ``$<$10\% dropout'' threshold without writing code. \textbf{ECG quality.} The Philips PageWriter~TC30 \cite{philips_tc30_manual} embeds automated 12-lead interpretation in each \texttt{.hea} file. The dashboard aggregates these verdicts and reports normal-vs-flagged proportions with HR and QTc distributions, so a researcher can decide whether to exclude, e.g., QTc~$>$~470\,ms or atrial fibrillation before training a sinus-rhythm-dependent model. Firmware (A.07.07.07) and filter settings were identical across all three sites, ruling out site-level hardware drift (Fig.~\ref{fig:signal_quality_module}). Both samples use the first participants in file ordering (N=200 CGM, N=150 ECG); full-cohort statistics require the offline batch pipeline noted in Section~\ref{sec:discussions}.

\begin{figure}[htbp]
    \centering
    \includegraphics[width=\linewidth]{Figure 4.png}
    \caption{Signal Quality module showing CGM dropout and ECG quality distributions with full sensor provenance for the Philips PageWriter TC30.}
    \label{fig:signal_quality_module}
\end{figure}

\section{Validation}
\label{sec:validation}
We assess validity along four axes rather than relying on interface appearance alone.

\textbf{(V1) Parser correctness.} Running the WFDB comment parser across all ECG participants (N=2,251 with confirmed \texttt{.hea} files) yielded a 100.0\% \texttt{validation\_date} extraction rate with zero failures across all failure modes (missing field, malformed date, header read error), confirming that the Philips TC30 embeds the comment consistently across all three sites. An offline batch script reproduces this figure over the full release and returns a per-failure-mode breakdown.

\textbf{(V2) Unit consistency.} As detailed in Section~\ref{sec:engineering}, the \texttt{unit\_source\_value} audit confirmed HbA1c and 9 further concepts unit-homogeneous, with non-canonical units isolated to four concepts and handled by range filtering; the cohort-level implausible-row exclusion rate is $<$0.2\%.

\textbf{(V3) Construct validity of labels.} The OMOP-derived HbA1c distribution shows the expected monotonic gradient across study groups (Section~\ref{sec:pipeline}, Step~5), and the insulin-dependent median of 7.0\% coincides with the ADA HbA1c treatment target evidence that the study-group labels are biomarker-consistent and usable as supervised targets.

\textbf{(V4) Cross-site consistency.} The site~$\times$~study-group breakdown (Table~\ref{tab:study_group}) shows uniformly high coverage (ECG 97.3--100.0\%, CGM 95.5--100.0\%, concurrent ECG+CGM 95.1--99.3\%), so no single site or group silently dominates a filtered cohort.

\section{Use Case: Preparing a Non-Invasive Glucose Prediction Pipeline}
\label{sec:pipeline}
We walk through a concrete scenario: building a training cohort for an ECG-to-glucose model. Table~\ref{tab:timesavings} summarises the per-task time, comparing the dashboard against an author-estimated manual baseline for a proficient Python/OMOP user; these estimates characterise the effort the tool removes and are not a substitute for the formal user study noted in Section~\ref{sec:discussions}.

\begin{table}[t]
    \centering
    \caption{Author-estimated task time: manual scripting vs.\ OpenT2D Explorer (reference workstation, warm cache).}
    \label{tab:timesavings}
    \begin{tabular}{@{} l c c @{}}
        \hline
        \textbf{Task} & \textbf{Manual (est.)} & \textbf{Dashboard} \\
        \hline
        Cohort scoping + coverage     & 2--4\,h        & $<$1\,min \\
        Modality triple-overlap       & 0.5--1\,day    & $<$1\,min \\
        Temporal check (per part.)    & 1--2\,h        & $\sim$30\,s \\
        Signal-quality gating         & 1--2\,days     & 9--18\,s \\
        Label validation (HbA1c)      & 1--2\,h        & $<$5\,s \\
        \hline
        \textbf{End-to-end}           & \textbf{several days} & \textbf{$\sim$15\,min} \\
        \hline
    \end{tabular}
\end{table}

\paragraph{Step~1 : Cohort scoping} The Overview page shows 2,280 participants across four study groups, with ECG and CGM each $>$95\% present; the researcher restricts initial experiments to the recommended \emph{train} split.

\paragraph{Step~2 : Modality intersection.} Setting the Co-missingness Matrix filter to the insulin-dependent group (highest severity, most clinically relevant) reveals that 248/258 (96.1\%) have concurrent ECG, CGM, and clinical files above the 200-sample minimum and the researcher proceeds.

\paragraph{Step~3 : Temporal co-registration.} Entering Participant~1023 in the TOV renders: visit 2023-08-30, CGM 2023-08-30 to 2023-09-09 (same calendar day, 9.9\,days, 1.4\% dropout, 2,816 readings), ECG 2024-10-14 (a $\approx$14-month follow-up). This confirms the usage pattern: CGM is consistently baseline-aligned, whereas ECG may be a separate longitudinal point and must be verified per participant.

\paragraph{Step~4 : Signal-quality gating.} The CGM histogram shows 99\% of the sample (198/200) with $<$10\% dropout (mean 0.4\%). The ECG panel reports the Philips four-way verdict distribution: 34.7\% \emph{Normal}, 26.0\% \emph{otherwise normal}, 6.0\% \emph{borderline}, and 33.3\% \emph{abnormal} (LVH, fascicular block, prolonged PR, old-infarct patterns consistent with a middle-aged T2D cohort), mean HR 68.3\,bpm, mean QTc 424.8\,ms. The researcher sets exclusion criteria (CGM dropout~$>$~20\%; strict \emph{abnormal} verdict with atrial fibrillation, or QTc~$>$~500\,ms) directly from the dashboard, without preprocessing code.

\paragraph{Step~5 : Label validation.} The Lab Distributions tab queries \texttt{measurement.csv} for HbA1c by group and shows a monotonic gradient: Healthy (median 5.5\%, IQR 5.3--5.7, n=759), Pre-DM (5.8\%, 5.6--6.0, n=546), Oral~Med.\ (6.3\%, 5.9--6.9, n=657), Insulin (7.0\%, 6.3--7.9, n=248). The insulin-dependent median matches the ADA target, confirming biomarker-validated labels.

This workflow cohort scoping, modality intersection, temporal verification, quality gating, and label validation takes $\approx$15\,min and zero custom code; the equivalent from raw files would take several days of parsing and exploration.

\section{Discussion}
\label{sec:discussions}
To our knowledge, OpenT2D Explorer is the first dedicated dashboard for the AI-READI dataset. Relative to the consortium Jupyter notebooks and OMOP-only tools such as ATLAS, its contribution is not the clinical numbers in Section~\ref{sec:pipeline} but the reusable three-axis method and the infrastructure that lets researchers without specialised parsing skills reach a quality-verified cohort.

\textbf{Maintainability, extensibility, and deployment.} The backend is organised as seven single-responsibility service modules behind three routers; adding a new modality means adding one service and one endpoint without touching existing parsers, and the OMOP conversion registry is designed for future mixed-unit releases. The entire system deploys via a single \texttt{docker compose up} with the dataset bind-mounted read-only, which both preserves DUA compliance and makes institutional deployment (shared VM or on-premises server) straightforward; it has been validated on a standard 4-vCPU cloud VM.

\textbf{Limitations.} The Co-missingness Matrix reflects file presence, not timestamp overlap, and TOV alignment is assessed only at calendar-date granularity: CGM timestamps are UTC (ISO~8601 with Z suffix), OMOP \texttt{visit\_start\_date} is timezone-naive, and the ECG \texttt{validation\_date} is a bare YYYYMMDD string. Fully resolving sub-day alignment would require site-specific UTC offsets and historical daylight-saving logs for UW/UCSD (Pacific) and UAB (Central), which are not in the v3.0.0 release; we treat calendar-date co-registration as the defensible operational threshold. The Signal Quality module samples the first 200 (CGM) and 150 (ECG) participants in file ordering, a head-slice chosen for interactive latency, not a random draw, so its rates are approximate indicators pending the offline full-cohort pipeline. Finally, the task-time analysis in Table~\ref{tab:timesavings} is an author benchmark; a formal usability study with clinical researchers (task-completion time, error rate, and satisfaction across multiple users) is the primary planned evaluation.

Table~\ref{tab:study_group} gives the site- and study-group-stratified ECG/CGM presence used in V4; coverage is uniformly high (lowest: UCSD Healthy CGM at 95.5\%), enabling site-level discrepancy checks before finalising a cohort. An automated flag for suspected timezone inconsistencies (ECG and CGM UTC calendar dates differing by exactly one day) is planned.

\begin{table}[htbp]
    \centering
    \caption{Site $\times$ Study-Group Stratified ECG and CGM File Presence}
    \label{tab:study_group}
    \resizebox{\columnwidth}{!}{%
    \begin{tabular}{@{} c c c c c c @{}}
    \hline
    \textbf{Site} & \textbf{Study Group} & \textbf{N} & \textbf{ECG (\%)} & \textbf{CGM (\%)} & \textbf{Both (\%)} \\
    \hline
    \multirow{4}{*}{UW (798)}   & Healthy    & 337 & 97.3 & 99.1 & 96.4 \\
                                 & Pre-DM     & 211 & 99.1 & 98.6 & 97.6 \\
                                 & Oral Med.  & 190 & 98.9 & 100.0 & 98.9 \\
                                 & Insulin    &  60 & 100.0 & 96.7 & 96.7 \\
    \hline
    \multirow{4}{*}{UCSD (682)} & Healthy    & 201 & 99.5 & 95.5 & 95.5 \\
                                 & Pre-DM     & 201 & 98.0 & 98.0 & 96.5 \\
                                 & Oral Med.  & 219 & 99.1 & 98.6 & 98.2 \\
                                 & Insulin    &  61 & 98.4 & 96.7 & 95.1 \\
    \hline
    \multirow{4}{*}{UAB (800)}  & Healthy    & 238 & 99.2 & 98.7 & 97.9 \\
                                 & Pre-DM     & 148 & 99.3 & 99.3 & 99.3 \\
                                 & Oral Med.  & 277 & 98.9 & 99.3 & 98.2 \\
                                 & Insulin    & 137 & 98.5 & 97.8 & 96.4 \\
    \hline
\end{tabular}
}
\end{table}

\section{Conclusion}
We presented OpenT2D Explorer, an open-source platform that lowers the barrier to data readiness in multimodal T2D studies. Its three EDA tools turn modality coexistence, temporal synchronisation, and signal integrity into a $\approx$15-minute, no-code process, and its WFDB timestamp-recovery technique makes cross-modal temporal alignment possible where standard headers fail. A cohort analysis showed 97.4\% triple-modality file presence and a clinically meaningful HbA1c gradient across study groups. Future work includes: (1)~timestamp-level temporal synchronisation for the full cohort; (2)~export to compute-oriented formats (Apache Parquet, WebDataset) consumable by PyTorch and TensorFlow pipelines; (3)~a cohort builder honouring the recommended train/val/test split; and (4)~a formal usability study with clinical researchers. The dashboard is available under the MIT License at \url{https://github.com/mdbasit897/aireadi-dashboard}.

\section*{Acknowledgment}
The authors thank the AI-READI Consortium (funded by NIH Bridge2AI) and the study participants. The AI-READI dataset is licensed under a customised FAIRhub license, and this dashboard is an analysis-only portal that does not republish any original data. This work is supported by the Google Cloud Research Credits program, award no.\ GCP19980904.

\begin{thebibliography}{00}

\bibitem{ai-readi_consortium_ai-readi_2024}
{AI-READI Consortium} \emph{et al.}, ``{AI}-{READI}: rethinking {AI} data collection, preparation and sharing in diabetes research and beyond,'' \emph{Nature Metabolism}, vol.~6, pp.~2210--2212, Nov.\ 2024. [Online]. Available: \url{https://doi.org/10.1038/s42255-024-01165-x}

\bibitem{ai-readi_consortium_flagship_2025}
{AI-READI Consortium}, ``Flagship Dataset of Type~2 Diabetes from the AI-READI Project,'' 2025. [Dataset]. Available: \url{https://fairhub.io/datasets/3/access}. Accessed: Mar~24, 2026.

\bibitem{wilkinson_fair_2016}
M.~D. Wilkinson \emph{et al.}, ``The FAIR Guiding Principles for scientific data management and stewardship,'' \emph{Scientific Data}, vol.~3, art.~160018, 2016. [Online]. Available: \url{https://doi.org/10.1038/sdata.2016.18}

\bibitem{moody_wfdb_nodate}
G.~Moody, T.~Pollard, and B.~Moody, ``{WFDB} Software Package,'' \emph{PhysioNet}, Jun.\ 2022, Version 10.7.0. [Online]. Available: \url{https://doi.org/10.13026/gjvw-1m31}

\bibitem{hripcsak_george_observational_2015}
G.~Hripcsak \emph{et al.}, ``Observational Health Data Sciences and Informatics (OHDSI): Opportunities for Observational Researchers,'' in \emph{Studies in Health Technology and Informatics}, vol.~216, pp.~574--578, 2015.

\bibitem{noauthor_omop_nodate}
OHDSI, ``OMOP Common Data Model v5.4.'' [Online]. Available: \url{https://ohdsi.github.io/CommonDataModel/cdm54.html}. Accessed: Apr~4, 2026.

\bibitem{ydata_profiling}
{YData}, ``ydata-profiling: Data profiling and exploratory data analysis in Python.'' [Online]. Available: \url{https://github.com/ydataai/ydata-profiling}. Accessed: Jul~18, 2026.

\bibitem{great_expectations}
{Great Expectations}, ``Great Expectations: Always know what to expect from your data.'' [Online]. Available: \url{https://greatexpectations.io/}. Accessed: Jul~18, 2026.

\bibitem{johnson_mimic_2023}
A.~E.~W. Johnson \emph{et al.}, ``MIMIC-IV, a freely accessible electronic health record dataset,'' \emph{Scientific Data}, vol.~10, art.~1, 2023. [Online]. Available: \url{https://doi.org/10.1038/s41597-022-01899-x}

\bibitem{pollard_eicu_2018}
T.~J. Pollard \emph{et al.}, ``The eICU Collaborative Research Database, a freely available multi-center database for critical care research,'' \emph{Scientific Data}, vol.~5, art.~180178, 2018. [Online]. Available: \url{https://doi.org/10.1038/sdata.2018.178}

\bibitem{noauthor_fastapi_nodate}
S.~Ram\'irez, ``FastAPI.'' [Online]. Available: \url{https://fastapi.tiangolo.com/}. Accessed: Mar~24, 2026.

\bibitem{dexcom_guides}
{Dexcom, Inc.}, ``Dexcom Product Guides.'' [Online]. Available: \url{https://www.dexcom.com/guides}. Accessed: Apr~4, 2026.

\bibitem{philips_tc30_manual}
{Philips Healthcare}, \emph{PageWriter TC30 Cardiograph --- User manual}. [Online]. Available: \url{http://www.frankshospitalworkshop.com/equipment/documents/ecg/user_manuals/Philips\%20PageWriter\%20TC30\%20Cardiograph\%20-\%20User\%20manual.pdf}. Accessed: Apr~4, 2026.

\end{thebibliography}

\end{document}
